\documentclass[11pt]{article}

\usepackage[a4paper,margin=29mm]{geometry}
\usepackage[T1]{fontenc}
\usepackage{lmodern}
\usepackage{microtype}
\usepackage{amsmath,amssymb,amsthm,mathtools}
\usepackage{array,booktabs}
\usepackage{enumitem}
\usepackage[hidelinks]{hyperref}


\newtheorem{theorem}{Theorem}[section]
\newtheorem{proposition}[theorem]{Proposition}
\newtheorem{lemma}[theorem]{Lemma}
\theoremstyle{definition}
\newtheorem{definition}[theorem]{Definition}
\theoremstyle{remark}
\newtheorem{remark}[theorem]{Remark}

\newcommand{\F}{\mathbb F}
\newcommand{\PG}{\operatorname{PG}}
\newcommand{\PGL}{\operatorname{PGL}}
\newcommand{\Aut}{\operatorname{Aut}}
\newcommand{\cC}{\mathcal C}
\newcommand{\supp}{\operatorname{supp}}

\title{Minimum-word reconstruction of $\PG(2,13)$\\
from a binary conic code}
\author{Tavis Rudd}
\date{August 2026}

\begin{document}
\begingroup
\makeatletter
\let\clebschmaintitle\@title
\def\@title{\large\textsc{The Clebsch cubic --- }\textbf{IV}\\[0.7em]
  \LARGE\clebschmaintitle}
\makeatother
\maketitle
\endgroup

\begin{center}
\small\itshape
From deep holes, the cubic takes shape, finds its bearings,\\
stands fixed while its shadows move, \textbf{rebuilds the plane that cast it},\\
and gathers its shadows home.
\end{center}

\begin{abstract}
Let \(\cC\) be a nonsingular conic in \(\PG(2,13)\), and form the binary
incidence code on its \(78\) internal points using the \(78\) passant lines.
We prove that the code has parameters \([78,36,12]_2\) and exactly \(364\)
minimum words.  Their weighted pair concurrences alone reconstruct the
passant incidence matrix, the code, and the six-class elliptic association
scheme.  The resulting group action then reconstructs all points and lines
of \(\PG(2,13)\), the distinguished conic, and its polarity; no coordinates
or triple concurrence are required.  Equivalently, the weighted \(2\)-section
of the minimum-support hypergraph is a complete invariant of this marked
conic-plane presentation.  The four minimum-word families are one
octahedral family and three chord-indexed punctured-conic families, and each
spans the code.  Structurally, the code is \(12\)-dimensional over a canonical
operator field \(\F_8\) --- a form of Madison and Wu's module decomposition ---
and that scalar action is why every family spans; we mark the scalar and build
the action from minimum-word pair data.  Exact positive-semidefinite and line-moment
certificates exclude weights eight and ten, replacing the corresponding
subset and syndrome searches.
\end{abstract}

\noindent\textbf{Keywords.} binary incidence code; finite conic; passant
line; minimum word; association scheme; reconstruction; module
decomposition; polarity.

\smallskip
\noindent\textbf{2020 Mathematics Subject Classification.} 94B05, 51E20,
05B25, 05E30.

\section{Introduction and statement}

Incidence codes usually remember less than the structures that define them.
Row operations erase the chosen parity checks, and minimum distance alone
does not explain the geometry of the first codewords.  The code studied here
behaves in the opposite way: the weighted pairs inside its minimum supports
recover not only the parity checks but the full projective plane in which
they arose.
This is the fourth paper in a program on code-theoretic reconstruction around
the Clebsch cubic.  The first paper begins from the conic code over
\(\F_{11}\); the present \(q=13\) theorem is logically independent of the
earlier papers.

Fix a nonsingular conic \(\cC\subset\PG(2,13)\).  A line is
\emph{passant} when it is disjoint from \(\cC(\F_{13})\).  Conic polarity
identifies the \(78\) internal points with the \(78\) passant lines.  Let
\(M\) be their binary incidence matrix and put
\[
 K=\ker_{\F_2}M\subseteq\F_2^{78}.
\]
The coordinate set is the set of internal points.  Automorphisms below are
coordinate permutations preserving the indicated code or minimum-support
system.

Let
\[
 \mathcal H=\{\supp(c):c\in K,\ \operatorname{wt}(c)=12\}
\]
be the minimum-support hypergraph on the unlabeled coordinate set of \(K\).

\begin{theorem}[Minimum words recover the conic plane]
\label{thm:main}
The code \(K\) has parameters \([78,36,12]_2\) and exactly \(364\)
minimum words.  Under \(\PGL(2,13)\), those words form four orbits of
size \(91\).  One is an octahedral homogeneous space with stabilizer
\(S_4\); the other three are chord-indexed punctured conics with stabilizer
\(D_{24}\).  Every orbit spans \(K\).

Let \(c(P,Q)\) be the number of minimum supports containing \(P,Q\).
The neighborhoods defined by \(c(P,Q)=8\) are exactly the \(78\) rows of
\(M\).  The other five elliptic relations are recovered from the remaining
pair weights and one common-neighbor count.  Moreover, if
\(P_{PQ}=c(P,Q)\bmod2\) off the diagonal and \(P_{PP}=0\), then
\(\operatorname{im}P=K\).  Thus the weighted \(2\)-section of
\(\mathcal H\) reconstructs the incidence matrix, code, and elliptic scheme.

From the reconstructed group, its fourteen Sylow-\(13\) subgroups and
\(169\) involutions recover all \(183\) points and lines of
\(\PG(2,13)\), the distinguished conic, and its polarity.  Consequently
weighted-pair reconstruction has exact arity two, and
\[
 \Aut(K)=\Aut(\mathcal H)=\PGL(2,13)
\]
in its symmetric-square action on the \(78\) internal points.
\end{theorem}

Reconstruction is intrinsic, not coordinatized.  It recovers the marked
plane and polarity up to isomorphism, but it does not choose an ordered
projective frame, a field labeling, or an equation for the conic.  Such a
choice is necessarily noncanonical: the \(2184\) ordered triples of distinct
conic points form a torsor for \(\PGL(2,13)\).

The proof is structural except at small exact endpoints.  A rank-\(28\)
positive-semidefinite matrix excludes weight eight.  A global line moment
reduces weight ten to two shapes, which are eliminated on four line-
stabilizer orbits and thirty-three point-stabilizer orbits.  The minimum
layer itself is exhausted at one point, the octahedral construction has one
bounded matching-profile check, and the recovered group/plane dictionary is
independently replayed in the complete finite model.

The code belongs to the fixed-conic LDPC family introduced by
Droms--Mellinger--Meyer~\cite{DromsMellingerMeyer2006}.  Their general bounds,
as recorded in the later comparison table of Ma--Liu--Tian
\cite[Table~1]{MaLiuTian2024}, specialize here to \(8\le d\le12\).
Madison and Wu proved the general nullity formula
\((q-1)^2/4\)~\cite[Corollary~6.3]{MadisonWu2012}; their Theorem~6.1(i) splits
the code over \(\overline{\F_2}\) into three non-isomorphic simple
twelve-dimensional \(\operatorname{PSL}(2,13)\)-modules.  These form one
Frobenius orbit, so the \(\F_2\)-form is irreducible of dimension thirty-six
with endomorphism field \(\F_8\), and spanning by any minimum-word family
follows.  We add the marking of \(A_9\) among the three conjugates and its
recovery from pair data.  Hollmann and Xiang constructed the elliptic
association scheme afforded by the same conic
action~\cite[Sections~3--5]{HollmannXiang2006}, attributing its first
description to Hollmann's thesis~\cite{HollmannThesis1982}.  The identification of
involutions with off-conic points, with involution axis equal to polar line,
is classical and is stated explicitly by Tranchida
\cite[Section~2]{Tranchida2024}.  We determine the upper endpoint of the
distance interval at \(q=13\), classify its minimum words, and prove that
their weighted pair data recovers these established ambient structures.

\section{Distance twelve}
\label{sec:distance}

We use the model \(\cC:XZ-Y^2=0\).  Every row and column of \(M\) has
weight seven, and conic polarity identifies the two index sets.  The
Madison--Wu formula gives \(\dim K=36\); direct binary elimination is an
independent check of rank \(42\).

Summing all row equations gives the all-one functional, so every word of
\(K\) has even weight.  If a nonzero word contains an internal point \(P\),
each of the seven passants through \(P\) contains a second support point.
Those seven points are distinct.  Thus every nonzero word has weight at
least eight.

\subsection{The tangent obstruction in weight eight}

Suppose that a word has weight eight and fix a support point \(P\).  Each of
the seven passants through \(P\) must contain another support point.  Since
only seven points remain, each passant contains exactly one of them.  Applying
the same argument at every support point shows that every support join is
passant and that no three support points are collinear.  Thus the support is
an eight-arc, and every passant pencil at a support point is saturated.  The
seven remaining lines through each support point are therefore both the
combinatorial tangents to the arc and the conic secants through that point.

For each secant line choose a nonzero linear equation.  For an internal point
\(P\), put
\[
 T_P(X)=\prod_{\substack{\ell\ni P\\
                    \ell\text{ secant to }\cC}}\ell(X).
\]
For a pairwise-passant triple define
\[
 h(P,Q,R)=
 \frac{T_P(Q)T_Q(R)T_R(P)}
      {T_P(R)T_R(Q)T_Q(P)}.
\]
Every displayed evaluation is nonzero because the three joins are passant.
Rescaling a line equation multiplies one numerator and one denominator factor
by the same scalar, so \(h\) is well defined.  The saturated-pencil
identification above makes the factors of \(T_P\) precisely the combinatorial
tangents at \(P\).
Segre's lemma of tangents in the coordinate-free form of
Ball--Lavrauw~\cite[Lemma~27]{BallLavrauw2019} gives
\(h(P,Q,R)=1\) for every triple in the hypothetical support.  Indeed, its
hypotheses apply to this eight-arc over the odd field \(\F_{13}\), and with
the cyclic order \(P,Q,R\) its tangent-product quotient is exactly the
displayed \(h\) (not its negative; reversing the cycle only inverts it).  Fix
\(P=(1:0:2)\).  Its other seven points would form a clique in the graph on
the \(42\) passant-join neighbors of \(P\), where \(Q\) and \(R\) are
adjacent when \(QR\) is passant and \(h(P,Q,R)=1\).

The projective map
\[
 \gamma(x:y:z)=
 (x+6y+9z:7x+9y+3z:10x+y+z)
\]
fixes \(P\) and has order \(14\) on its passant-join neighbors.  With indices
in \(\mathbf Z/14\), put
\[
 A_i=\gamma^i(1:1:7),\qquad
 B_i=\gamma^i(1:1:6),\qquad
 C_i=\gamma^i(1:2:10).
\]
These three orbits are disjoint and contain all \(42\) neighbors.  Direct
substitution in the product above gives
\[
\begin{array}{c|c}
\text{orbit pair}&j-i\pmod {14}\\ \hline
AA&4,6,8,10\\
AB&6,7,11,12\\
AC&1,3\\
BB&6,8\\
BC&3,5,6,8,9,11\\
CC&2,4,10,12.
\end{array}
\]
Let \(C\) be the symmetric \(42\)-square matrix in the displayed vertex
order, written in \(14\)-square circulant blocks.  Its six upper-block first
rows are
\[
\begin{array}{c|rrrrrrrrrrrrrr}
00&40&9&8&-8&-10&0&-10&-2&-10&0&-10&-8&8&9\\
01&7&6&6&6&7&-1&-10&-10&-8&0&-8&-10&-10&-1\\
02&0&-10&2&-10&0&3&2&9&2&4&2&9&2&3\\
11&40&15&10&10&-6&-3&-10&-22&-10&-3&-6&10&10&15\\
12&-13&3&0&-10&8&-10&-10&14&-10&-10&8&-10&0&3\\
22&40&-9&-10&19&-10&2&8&-12&8&2&-10&19&-10&-9.
\end{array}
\]
The lower blocks are the transposes.  Comparison with the six difference
sets shows that \(C\) has diagonal \(40\) and entry \(-10\) on every edge
of the local graph.

The matrix is positive semidefinite.  Exact cyclic Fourier traces and Newton
identities give
\[
 \chi_C(x)=x^{14}(x^2-94x+279)(x^2-26x+135)f(x)^2g(x)^2,
\]
where
\[
\begin{aligned}
f(x)={}&x^6-247x^5+22466x^4-910807x^3
 +15948097x^2-111043270x+189747571,\\
g(x)={}&x^6-533x^5+105578x^4-9736765x^3
 +438572569x^2-9062718662x+68467048091.
\end{aligned}
\]
Every nonzero factor has alternating coefficients.  Its value at \(-t\) is
therefore positive for \(t\ge0\).  Since \(C\) is real symmetric, it has no
negative eigenvalue; the displayed factorization also gives rank \(28\).
An independent rational Schur elimination gives the same twenty-eight
positive pivots.

If \(v\) is the characteristic vector of a clique of size \(c>0\), then
\[
 0\le v^{\mathsf T}Cv
   =40c-20\binom c2=10c(5-c).
\]
Thus \(c\le5\).  The set
\(\{A_0,B_6,B_{12},C_1,C_3\}\) is a five-clique, so the bound is sharp.
Moreover, its fourteen translates are linearly independent kernel vectors;
nullity fourteen makes them a basis of \(\ker C\).  Equality in the
quadratic-form bound puts every maximum-clique vector in this kernel, and
solving the resulting \(0\)--\(1\) constraints in that basis gives precisely
the fourteen translates.  In particular the seven-clique required by a
weight-eight word cannot occur.

\subsection{The two weight-ten profiles}

Let \(S\) be a hypothetical weight-ten support, and let \(n_i\) be the
number of passants meeting \(S\) in exactly \(i\) points.  Code parity makes
every occurring \(i\) even, and incidence counting gives
\[
 \sum_i i n_i=70.
\]
At a selected point, pencil parity gives secant degree zero or two.  The
secant graph induced on \(S\) is therefore a union of cycles and isolated
vertices.  If \(m\) is the number of its cycle vertices, then it has \(m\)
edges, and the other \(45-m\) pairs have passant join.  Hence
\[
 \sum_i\binom{i}{2}n_i=45-m.
\]
Subtracting half the first identity yields the decisive moment
\begin{equation}
 4n_4+12n_6+24n_8+\cdots=10-m.       \label{eq:weight-ten-moment}
\end{equation}
No \(n_i\) with \(i\ge6\) can occur.  The degree condition and
\eqref{eq:weight-ten-moment} leave exactly
\[
\begin{array}{c|c|c}
m&(n_2,n_4)&\text{shape}\\ \hline
6&(33,1)&\text{four isolated vertices on one four-point passant and six cycle vertices},\\
10&(35,0)&\text{ten cycle vertices; every passant meets }S\text{ in }0\text{ or }2.
\end{array}
\]
Indeed two degree-two vertices cannot by themselves form a cycle; in the
first case uniqueness of the four-point passant puts all isolated vertices
on that line.

The first shape has only four normalized cases.  The effective stabilizer of
a passant acts as \(D_{14}\) on its seven internal points, with four orbits
on four-subsets:
\[
\begin{array}{c|c|c|c}
\text{representative}&\text{orbit size}&|X|&
 \text{secant degrees in }X\\ \hline
0123&7&7&(3,3,4,4,4,4,4)\\
0124&14&3&(0,0,0)\\
0134&7&5&(2,2,2,4,4)\\
0135&7&5&(2,2,2,3,3).
\end{array}
\]
Here \(X\) is the set of internal points passantly joined to all four fixed
points; retaining the three unused points of the fixed line only enlarges
it.  None of the four induced secant graphs contains a six-vertex
two-regular subgraph, so the \(m=6\) shape is impossible.

For \(m=10\), fix \(P\in S\).  The remaining points consist of one point in
each of the seven six-point passant fibres through \(P\) and an unordered
pair among its \(35\) secant neighbors.  The stabilizer
\(G_P\cong D_{28}\) has thirty-three orbits on the \(595\) pairs: five of
size \(7\), sixteen of size \(14\), and twelve of size \(28\).  For one
representative of each orbit, add the seven fibre points in order, rejecting
a prefix once a passant contains three selected points or a selected point
has more than two secant neighbors.  The complete certificate has last
surviving-depth mass
\[
\begin{array}{c|rrrrr}
\text{depth}&0&1&2&3&4\\ \hline
\text{pair mass}&28&35&91&238&203.
\end{array}
\]
The masses sum to \(595\), and no representative survives the fifth fibre.
Thus \(m=10\) is impossible as well.

Finally, the twelve internal points
\[
\begin{gathered}
(1:0:2),(1:3:2),(1:4:5),(1:1:8),
(1:4:8),(1:1:7),\\
(1:7:12),(1:3:3),(1:9:11),(1:10:11),
(1:0:5),(1:8:7)
\end{gathered}
\]
have zero syndrome.  Their stabilizer is dihedral of order \(24\).
Therefore \(d(K)=12\).

\section{Minimum geometry and weighted pairs}

Fix one support point \(P\).  Its seven passant fibres again have six eligible
points each, and it has \(35\) secant neighbors.  Pencil parity leaves
\(s=0,2,4\) secant neighbors.
The corresponding multiplicities in the seven passant fibres, followed by
the number of secant neighbors, are
\[
 (5,1^6;0),\qquad (3,3,1^5;0),\qquad
 (3,1^6;2),\qquad (1^7;4).
\]
A syndrome-indexed exact check exhausts these four fixed-point domains.  Its
input sizes are respectively
\[
 7\binom65 6^6,\quad
 \binom72\binom63^2 6^5,\quad
 7\binom63 6^6\binom{35}{2},\quad
 6^7\binom{35}{2}^2,
\]
with disjoint secant pairs in the last domain.  The certificate specifies the
four partitions, proves their coverage, tests zero syndrome, and deduplicates
only after conversion to unordered supports.  The four domains contribute
\(0,0,0,56\) words.  The canonical replay generates the four projective
\(\PGL(2,13)\)-orbits, and intersects each orbit with the supports containing
\(P\).  These four fixed-point slices are disjoint, each has size \(14\), and
their union is checked for equality with the exhaustive \(56\)-word set.
The order-\(28\) stabilizer of the fixed point acts transitively on each
slice; the stabilizer of a support inside it has order \(2\).
Transitivity, or equivalently the count \(56\cdot78/12=364\), proves global
exhaustion.

The replay generates each orbit from its recorded representative and checks
every support against \(M\).  The following constructions identify those
orbits without representatives.

\subsection{One octahedral and three toric families}

Parametrize the conic by \((t^2:t:1)\), with its usual point at infinity.
An internal point \(P=(x:y:z)\) defines the fixed-point-free involution
\[
 j_P(t)=\frac{yt-x}{zt-y},\qquad
 J_P=\begin{pmatrix}y&-x\\z&-y\end{pmatrix},\qquad
 J_P^2=(y^2-xz)I.
\]
The seven secants through \(P\) are the seven pairs
\(\{t,j_P(t)\}\).  Thus internal points may be read as perfect matchings of
the fourteen conic points.

Fix a chord \(e\), carry its endpoints to \(\{0,\infty\}\), and let
\(N_e\cong D_{24}\) be their setwise stabilizer.  The three supports over
\(e\) are
\[
 S_r(e)=\{(x:1:r^{-1}x^{-1}):x\in\F_{13}^{\times}\},
 \qquad r\in\{2,5,11\}.
\]
Intrinsically, these three parameters satisfy
\(\chi(r)=-1\) and \(\chi(r-1)=1\).  The support \(S_r(e)\) is the
punctured pencil conic
\[
 \{y^2=rxz\}(\F_{13})\setminus e,
\]
so it has size twelve and stabilizer exactly \(N_e\).  At each displayed
point, \(\Delta=1-r^{-1}=(r-1)/r\) is a nonsquare; hence all twelve points
are internal.  The set is a codeword for a structural reason.  If a passant
\(AX+BY+CZ=0\) met \(S_r(e)\) with odd multiplicity, the polynomial
\(Ax^2+Bx+C/r\) would have a double root.  Then
\[
 B^2-4AC=B^2(1-r)
\]
would be a square, contrary to the line being passant.  Hence every passant
meets \(S_r(e)\) in zero or two points.

For the remaining family, take the octahedral subgroup
\(O\cong S_4\) whose conic-point orbits have sizes six and eight.  Let
\(X_O\) consist of the internal points whose matching has two pairs inside
the six-set, three inside the eight-set, and two crossing pairs.  In a
standard octahedral frame, the six internal-point suborbits have sizes
\(6,12,12,12,12,24\); the displayed matching profile selects one of the
twelve-orbits, and direct incidence parity makes it a codeword.  Thus
\(|X_O|=12\) and \(\operatorname{Stab}_G(X_O)=O\).

All four stabilizers have order \(24\).  Each family therefore has
\(2184/24=91=\binom{14}{2}\) supports, agreeing with the fixed-point
exhaustion above.

\subsection{Exact arity-two reconstruction}

For distinct coordinates put
\[
 c(P,Q)=\#\{S\in\mathcal H:\{P,Q\}\subseteq S\}.
\]
The exact pair table is
\[
\begin{array}{c|rrrrrr}
\rho&0&1&3&9&10&12\\ \hline
c&8&6&6&12&7&9.
\end{array}
\]
Only the relations \(A_1\) and \(A_3\) are fused.  For a concurrence-six
pair define
\[
 d_7(P,Q)=\#\{R:c(P,R)=c(R,Q)=7\}.
\]
The \(A_{10}^2\) intersection row gives
\[
 d_7(P,Q)=2\quad\text{on }A_1,
 \qquad d_7(P,Q)=4\quad\text{on }A_3.
\]
This pair-derived common-neighbor count recovers the full six-color elliptic
scheme.

The incidence rows require no coherent closure.  Polarity sends
\(P=(x:y:z)\) to \(P^\perp=(-z:2y:-x)\), and
\[
 Q\in P^\perp
 \iff \beta(P,Q)=0
 \iff \rho(P,Q)=0
 \iff c(P,Q)=8.
\]
Consequently the seventy-eight concurrence-eight neighborhoods are exactly
the seventy-eight rows of \(M\).

The pair data also recovers the code without first naming a row.  Let
\(P_{PQ}=c(P,Q)\bmod2\) off the diagonal and \(P_{PP}=0\).  Then
\[
 P=A_{10}+A_{12},\qquad \operatorname{rank}P=36,
 \qquad A_0P=0,
\]
so \(\operatorname{im}P=K\).  Every coordinate occurs in
\(364\cdot12/78=56\) minimum supports, making unary data constant.  Pair
data suffices and unary data does not; reconstruction arity is exactly two.

\section{The hidden field, spanning, and automorphisms}

For internal points \(P=(x_P:y_P:z_P)\) and \(Q\), set
\[
 \Delta(P)=y_P^2-x_Pz_P,\qquad
 \beta(P,Q)=2y_Py_Q-x_Pz_Q-z_Px_Q,
\]
and
\[
 \rho(P,Q)=\frac{\beta(P,Q)^2}{\Delta(P)\Delta(Q)}.
\]
On distinct internal points, \(\rho\) takes the six values
\(0,1,3,9,10,12\).  Its level relations are the elliptic scheme; write
\(A_r\) for their adjacency matrices.  Polarity identifies \(A_0\) with
\(M\), up to row order.

The needed part of the integral intersection table reduces modulo two to
\begin{equation}
\label{eq:scheme}
\begin{split}
 A_0^2&=I+A_9+A_{10}+A_{12},\\
 A_0A_r&=0\quad(r=9,10,12),\\
 A_9^2&=A_{10},\qquad A_{10}^2=A_{12},\qquad A_{12}^2=A_9.
\end{split}
\end{equation}
These are parity reductions of intersection numbers.  The
\(\PGL(2,13)\)-action is transitive on each relation, so fixing one point and
substituting one representative of every relation checks all entries.

Put \(B=A_9\).  Since \(A_0B=0\), the image of \(B\) lies in \(K\).  If
\(x\in K\) and \(Bx=0\), then \eqref{eq:scheme} gives
\[
 0=A_0^2x=(I+B+B^2+B^4)x=x.
\]
Thus \(B\) is injective on the 36-dimensional space \(K\), so
\(\operatorname{im}B=K\), and \(B,B^2,B^4\) all have rank 36.

\begin{proposition}[The operator field]
\label{prop:hidden-field}
The action of the binary Bose--Mesner algebra on \(K\) has image \(\F_8\).
If \(\alpha\) is a root of \(t^3+t^2+1\), then
\[
 K\cong\F_8^{12},\qquad
 (A_9,A_{10},A_{12})=(\alpha,\alpha^2,\alpha^4).
\]
The identification of a named \(\alpha\) is canonical up to Frobenius.
\end{proposition}

Indeed, exact elimination gives \(\operatorname{rank}(B+I)=78\).  Restricting
the first identity in \eqref{eq:scheme} to \(K\) and factoring
\[
 1+t+t^2+t^4=(t+1)(t^3+t^2+1)
\]
therefore gives \(B^3+B^2+I=0\).  The cubic has no root in \(\F_2\), so it
is irreducible and makes \(K\) a vector space of dimension \(36/3=12\) over
\(\F_8\).  The squaring identities in \eqref{eq:scheme} give the three
displayed scalar operators.  This is an operator-field structure on the
binary code, not a relabeling as a length-26 code over \(\F_8\).

If \(N_i\) is the 91-by-78 support matrix of the \(i\)-th minimum-word
orbit, direct pair counting in one intrinsic representative gives
\[
 (N_1^{\mathsf T}N_1,N_2^{\mathsf T}N_2,
   N_3^{\mathsf T}N_3,N_4^{\mathsf T}N_4)
 =(A_9,A_9,A_{12},A_{10}).
\]
These are nonzero scalars on \(\F_8^{12}\), hence automorphisms of \(K\).
Thus \(\operatorname{im}(N_i^{\mathsf T}N_i)=K\) lies in the row space of
\(N_i\).  Its rows already lie in \(K\), so every orbit spans \(K\).

It remains to determine the scheme automorphisms.  Use the four anchors
\[
 P_0=(1:0:2),\quad P_1=(1:1:7),\quad
 P_2=(1:0:7),\quad P_3=(1:1:3).
\]
The rigidity mechanism is summarized by
\[
 \underbrace{(P_0,P_1,P_2)}_{(10,3,9)\text{ regular orbit}}
 \quad\Longrightarrow\quad
 \underbrace{P_3}_{(3,1,9)\text{ is unique}}
 \quad\Longrightarrow\quad
 \underbrace{P\longmapsto
   \bigl([P=P_i],\rho(P,P_i)\bigr)_{i=0}^{3}}
   _{\text{injective on the 78 points}}.
\]
The first three satisfy
\[
 \bigl(\rho(P_0,P_1),\rho(P_0,P_2),\rho(P_1,P_2)\bigr)=(10,3,9).
\]
There are
\(78\cdot14\cdot2=2184\) ordered triples with this pattern: \(A_{10}\)
has valency 14 and \(p_{3,9}^{10}=2\).  Substitution in the symmetric-square
action shows that the stabilizer of \((P_0,P_1,P_2)\) in
\(\PGL(2,13)\) is trivial.  Since the group also has order 2184, it acts
simply transitively on these triples.

For a fixed triple, \(P_3\) is the unique point with relation signature
\((3,1,9)\).  The four values
\[
 \bigl(\rho(P,P_0),\rho(P,P_1),\rho(P,P_2),\rho(P,P_3)\bigr),
\]
with equality to an anchor recorded on the diagonal, distinguish all 78
points.  Consequently a scheme automorphism can be composed with a unique
element of \(\PGL(2,13)\) to fix the first three anchors; it then fixes the
fourth and every point.  Thus the reconstructed scheme has automorphism
group \(\PGL(2,13)\).

Every coordinate automorphism of \(K\) preserves \(\mathcal H\).
Conversely, \(\mathcal H\) spans \(K\), so every automorphism of
\(\mathcal H\) preserves \(K\).  The two coordinate-permutation groups
coincide.

\section{Recovery of the ambient conic plane}

It remains to recover the plane rather than only its internal-point scheme.
Let
\[
 G=\Aut(\mathcal H)\cong\PGL(2,13),\qquad
 \Omega=\operatorname{Syl}_{13}(G).
\]
A Sylow-\(13\) subgroup has normalizer of order \(156\), so \(|\Omega|=14\).
Conjugation gives a sharply three-transitive action: in the standard model a
Sylow subgroup is the unipotent group fixing one point of
\(\mathbf P^1(\F_{13})\), and both \(G\) and the set of ordered distinct
triples have size \(2184\).  Thus \(\Omega\) is intrinsically the recovered
conic-point set.

Let \(\mathcal I\) be the set of \(169\) involutions of \(G\), and put
\(\Pi=\Omega\sqcup\mathcal I\).  Index both points and polar lines by
\(\Pi\).  Define incidence by
\begin{enumerate}[label=(\roman*)]
\item \(U\in U^\perp\) for \(U\in\Omega\), with no other
  \(\Omega\)--\(\Omega\) incidences;
\item \(U\in j^\perp\) for \(U\in\Omega\), \(j\in\mathcal I\), exactly
  when \(j\) normalizes \(U\);
\item \(i\in j^\perp\) for distinct \(i,j\in\mathcal I\) exactly when
  \(ij\) is an involution.
\end{enumerate}

These rules are the trace polarity on the adjoint module.  Identify the
standard plane with \(\mathbf P(\mathfrak{sl}_2(\F_{13}))\).  For traceless
\(A\), Cayley--Hamilton gives \(A^2=-\det(A)I\).  The fourteen singular
matrix lines are the nilpotent conic and correspond to the Sylow subgroups;
the nonsingular lines correspond bijectively to projective involutions.
Moreover
\[
 \langle A,B\rangle=\operatorname{tr}(AB)
\]
is the polar form of the determinant conic.  Orthogonality of two nilpotent
lines means equality, orthogonality of a nilpotent and an involution means
normalization of the Sylow subgroup, and for distinct involutions it is
equivalent to their product being an involution.  The three intrinsic rules
therefore recover \(\PG(2,13)\), its conic \(\Omega\), and its polarity.

The \(78\) original coordinates are also intrinsic inside this plane.  A
coordinate stabilizer is a nonsplit-torus normalizer \(D_{28}\); its unique
central involution identifies the coordinate with the involution class of
centralizer order \(28\).  The other involution class has size \(91\).
Consequently the old matrix \(M\) is exactly the internal--internal block of
the recovered polarity matrix.  Since Section~3 reconstructed \(G\) and
\(M\) from weighted pairs alone, the whole marked plane is reconstructed at
arity two.  This completes the proof of Theorem~\ref{thm:main}.

\section{Proof and verification boundary}

The human proof owns the reductions: tangent products produce the local
weight-eight graph, line moments produce the two weight-ten shapes, conic
pencils produce the toric words, weighted pairs produce the polarity rows,
and the adjoint representation produces the ambient plane.  Exact execution
enters only after those reductions.

The accompanying \nolinkurl{verification/} directory contains separate
canonical certificates for the theta matrix, moment--stabilizer leaves,
weighted-pair reconstruction, toric--octahedral geometry, ambient plane, and
hidden operator field.  The earlier clique and syndrome programs remain as
independent checks but are not theorem-facing proof leaves.  The minimum-word
exhaustion and four-anchor rigidity retain their original exact replays.
The entry point \nolinkurl{verify_evidence.py} checks byte counts and SHA-256
hashes before running every program.  All arithmetic is exact and
deterministic.
From the paper root, the public command
\begin{center}
\small\texttt{make check}
\end{center}
ends with \texttt{Paper-IV q=13 evidence: PASS} and then builds the
warning-free PDF.  The exact programs are described in
\nolinkurl{verification/README.md}; the formal package is distributed
separately and carries its own README.  The Python evidence uses only the
standard library, the manuscript build uses a Nix-pinned \TeX\ environment,
and the formal package a Nix-pinned Lean environment.

The shared Lean library formalizes the normalized conic and incidence code,
the theta inequality, moment compression, pair-neighborhood interface,
pencil discriminant identity, and hidden-cubic cancellation.  The
paper-owned package checks pair-only recovery, unary constancy, the three
toric parities, the hidden cubic on the relation-operator image, and both
projective-plane uniqueness axioms on the normalized \(183\)-point model.
Table~\ref{tab:trust} records the remaining division of proof.

\begin{table}[ht]
\centering
\footnotesize
\begin{tabular}{@{}>{\raggedright\arraybackslash}p{0.17\textwidth}
  >{\raggedright\arraybackslash}p{0.24\textwidth}
  >{\raggedright\arraybackslash}p{0.49\textwidth}@{}}
\toprule
claim & proof mode & exact boundary\\
\midrule
dimension \(36\)
& published theorem; Lean theorem
& Madison--Wu gives the formula; exact elimination and semantic recovery maps
  prove rank \(42\) at \(q=13\).\\
weight \(8\)
& human/classical reduction; exact PSD certificate; Lean implication
& Segre's lemma gives the 42-vertex graph.  Exact Fourier traces and rational
  Schur elimination certify the rank-28 form; Lean proves the quadratic-form
  clique implication.\\
weight \(10\)
& human moment proof; exact stabilizer leaves; Lean algebra
& Incidence and pair counts give \eqref{eq:weight-ten-moment}; four
  \(D_{14}\) and thirty-three \(D_{28}\) representatives close the two
  shapes.\\
minimum layer
& Lean native evaluation; trusted Python; human transport
& Both executions exhaust the \(56\) fixed-point words; projective
  transitivity and the double count give all \(364\).\\
minimum geometry
& human conic proof; bounded exact leaf; Lean parity checks
& The toric families have a discriminant proof.  One octahedral matching
  profile and the four stabilizers are checked in standard form.\\
pair reconstruction
& human polarity proof; trusted exact data; Lean theorem
& Pair counts and one common-neighbor count recover the scheme; color-eight
  neighborhoods recover the rows, and pair parity recovers the code.\\
ambient plane
& human adjoint proof; trusted exact replay; Lean incidence gate
& The Sylow/involution rules are identified with trace polarity.  Lean checks
  the 183-point plane axioms; the group-theoretic census is exact Python.\\
hidden field and spanning
& human algebra; trusted matrix data; Lean theorem
& The irreducible cubic makes \(K\cong\F_8^{12}\); checked Gram scalars and a
  kernel-checked image argument prove that every orbit spans.\\
\bottomrule
\end{tabular}
\caption{Proof modes and trust boundaries for Theorem~\ref{thm:main}.}
\label{tab:trust}
\end{table}

The paper-owned aggregate namespace is
\nolinkurl{PassantCodeQ13.Gates.Main}; its terminals are
\begin{flushleft}
\small
\nolinkurl{incidenceRankAndCodeDimension}\\
\nolinkurl{fixedPointWeightTwelveExhaustion}\\
\nolinkurl{minimumOrbitsSpanRhoZeroKernel}\\
\nolinkurl{ellipticSchemeAutomorphismsAreProjective}\\
\nolinkurl{pairOnlyReconstruction}\\
\nolinkurl{toricMinimumSupports}\\
\nolinkurl{hiddenFieldCubicOnImage}\\
\nolinkurl{ambientPlaneIncidence}.
\end{flushleft}
The shared low-weight namespace is
\nolinkurl{RelativeConicArcs.Gates.PassantCodeQ13}; its terminals are
\begin{flushleft}
\small
\nolinkurl{weightEight_semantic_transport}\\
\nolinkurl{arbitrary_weightTen_profile_transport}.
\end{flushleft}

Native evaluation compiles and runs finite decision procedures rather than
reducing them entirely in the kernel.  Under Lean \(4.32.0\)-rc1, the
tracked \texttt{\#print axioms} audit therefore exposes each such use as a
declaration-local axiom; the symbolic transport theorems remain
kernel-checked relative to those finite leaves.
Thus ``Lean theorem'' in Table~\ref{tab:trust} means kernel-checked relative
to its printed axioms, whereas ``Lean native evaluation'' and ``trusted
Python'' are executions.  Hash checks identify the executed sources and
outputs but do not turn either execution into a kernel-checked certificate;
no finite leaf in the table is described that way.
The formal scope retains three explicit human transports.  Projective
transitivity and double counting carry the fixed-point exhaustion to all
\(364\) minimum words.  Weighted-pair equivariance carries support
automorphisms to the polar-relation scheme.  The adjoint trace model
identifies the abstract Sylow/involution incidence with the conic polarity.
Segre's lemma of tangents remains the cited classical input to the
weight-eight transport.

\section{Conclusion}

The minimum words retain the geometry that row reduction forgets.  Their
weighted pairs recover the parity checks and code, then the elliptic scheme,
and finally the full conic plane and polarity.  The theta form and global
line moment explain why this reconstruction first appears in weight twelve;
the toric--octahedral families explain what appears there; and the hidden
\(\F_8\) action explains why each family spans.  Thus, in this \(q=13\)
instance, the weighted minimum-support \(2\)-section, passant incidence
matrix, binary code, elliptic scheme, and marked conic plane are mutually
recoverable up to the stated isomorphisms.

Within the Clebsch cubic program, this is the endpoint of the reconstruction
theme: a code-theoretic shadow recovers its complete ambient geometry.  The
argument is nevertheless standalone.  It uses neither the earlier papers nor
a chosen Clebsch model, and it claims no uniform minimum-distance theorem
beyond \(q=13\).

\begin{thebibliography}{9}

\bibitem{BallLavrauw2019}
S.~Ball and M.~Lavrauw,
\newblock Arcs in finite projective spaces,
\newblock \emph{EMS Surveys in Mathematical Sciences} \textbf{6} (2019),
133--172.
\newblock arXiv:1908.10772.

\bibitem{DromsMellingerMeyer2006}
S.~V.~Droms, K.~E.~Mellinger, and C.~Meyer,
\newblock LDPC codes generated by conics in the classical projective plane,
\newblock \emph{Designs, Codes and Cryptography} \textbf{40} (2006),
343--356.
\newblock doi:10.1007/s10623-006-0022-6.

\bibitem{HollmannThesis1982}
H.~D.~L.~Hollmann,
\newblock \emph{Association schemes},
\newblock Master's thesis, Eindhoven University of Technology (1982).

\bibitem{HollmannXiang2006}
H.~D.~L.~Hollmann and Q.~Xiang,
\newblock Association schemes from the action of
\(\PGL(2,q)\) fixing a nonsingular conic in \(\PG(2,q)\),
\newblock \emph{Journal of Algebraic Combinatorics} \textbf{24} (2006),
157--193.
\newblock doi:10.1007/s10801-006-0005-8; arXiv:math/0503573.

\bibitem{MadisonWu2012}
A.~L.~Madison and J.~Wu,
\newblock On binary codes from conics in \(\PG(2,q)\),
\newblock \emph{European Journal of Combinatorics} \textbf{33} (2012),
33--48.
\newblock doi:10.1016/j.ejc.2011.08.001; arXiv:1104.0324.

\bibitem{MaLiuTian2024}
L.~Ma, S.~Liu, and Z.~Tian,
\newblock The binary codes generated from quadrics in projective spaces,
\newblock \emph{AIMS Mathematics} \textbf{9} (2024), 29333--29345.
\newblock doi:10.3934/math.20241421.

\bibitem{Tranchida2024}
P.~Tranchida,
\newblock Triples of involutions in \(\PGL(2,q)\) and their incidence
geometries,
\newblock arXiv:2411.10299, 2024.

\end{thebibliography}

\enlargethispage{5\baselineskip}
\section*{AI assistance disclosure}

This project was developed with extensive assistance from OpenAI Codex and
Anthropic Claude. Under the author's direction, the systems assisted with
proof exploration, literature research, symbolic and finite computation,
code and formal-proof development, verification, and manuscript drafting and
revision. The author checked the resulting arguments, computations, code,
and cited sources, reviewed and edited all AI-assisted material, and assumes
responsibility for all content.

\end{document}
